\section[Estimates of the even moments via $U_n$]{Estimates of the even moments via $\boldsymbol{U_n}$}\label{section3}
In all of this section we assume that the moment problem is symmetric or equivalently that $a_n\equiv 0$, so the recurrence relation~\eqref{eq:rec} takes the simplified form
\begin{gather}\label{eq:recsym}
xP_n=b_nP_{n+1}+b_{n-1}P_{n-1}.
\end{gather}

\begin{prop}\label{thm:increasofbn}Let $(b_n)$ and $\big(\tilde{b}_n\big)$ denote two sequences of positive numbers satisfying $b_n\le \tilde{b}_n$, $n\ge 0$. Then the corresponding even moments satisfy $s_{2n}\le \tilde{s}_{2n}$, $n\ge 0$.
\end{prop}

\begin{proof}Let $J$ denote the Jacobi matrix associated to the moment problem, i.e.,
\begin{gather}\label{eq:Jacobimat}
J=\begin{pmatrix} 0 & b_0 & 0 & 0 & \cdots\\
 b_0 & 0 & b_1 & 0 & \cdots\\
 0 & b_1 & 0 & b_2 & \cdots\\
 0 & 0 & b_2 & 0 & \cdots\\
 \vdots & \vdots & \vdots & \vdots & \ddots
\end{pmatrix}.
\end{gather}
It is well-known that $s_n=(J^n\delta_0,\delta_0)$, and the conclusion follows.
\end{proof}

In the symmetric case $P_k$ involves only monomials with the same parity as $k$. Thus we have
 \begin{gather}\label{eq:u0}
{x^n\over b_0b_1\cdots b_{n-1}}=\sum_{k=0}^{[n/2]}u_{n,k}P_{n-2k}=u_{n,0}P_n+u_{n,1}P_{n-2}+\cdots,
\end{gather}
 hence
 \begin{gather}\label{eq:u1}
U_n:= {s_{2n}\over b_0^2b_1^2\cdots b_{n-1}^2}=\sum_{k=0}^{[n/2]}u_{n,k}^2.
 \end{gather}
(Compared with \eqref{eq:x} we have $u_{n,k}=c_{n-2k,n}/(b_0\cdots b_{n-1})$, while $c_{n-2k-1,n}=0$.) By \eqref{eq:lead} we have $u_{n,0}=1$.

By the recurrence relation \eqref{eq:recsym} the coefficients $u_{n,k}$ satisfy
\begin{gather}\label{eq:u}
u_{n+1,k}={b_{n-2k}\over b_n} u_{n,k}+{b_{n+1-2k}\over b_n}u_{n,k-1},\qquad 1\le k\le [n/2].
\end{gather}
In fact, multiplying \eqref{eq:u0} by $x/b_n$ and using the recurrence relation \eqref{eq:recsym}, we get
\begin{gather*}
\frac{x^{n+1}}{b_0b_1\cdots b_n}=\sum_{k=0}^{[n/2]}\frac{u_{n,k}}{b_n}\left(b_{n-2k}P_{n-2k+1}+b_{n-2k-1}P_{n-2k-1}\right),
\end{gather*}
and \eqref{eq:u} follows. We may extend the definition by setting $u_{n,k}=0$ for $k>[n/2]$ or $k<0$ and $b_n=0$ for $n<0$. In this way it is not difficult to see that
 \begin{gather}\label{eq:u'}
u_{n+1,k}={b_{n-2k}\over b_n} u_{n,k}+{b_{n+1-2k}\over b_n}u_{n,k-1},\qquad n, k\ge 0,
\end{gather}
and it follows by induction that all coefficients $u_{n,k}\ge 0$.

By \eqref{eq:u1} and \eqref{eq:u} this immediately implies the following.

\begin{prop}\label{thm:comparison} Assume $(b_n)$ and $\big(\tilde{b}_n\big)$ are positive sequences satisfying
\begin{gather*}
{b_{n-1}\over b_{n}} \le {\tilde{b}_{n-1}\over \tilde{b}_{n}} ,\qquad n\ge 1,
\end{gather*}
 and let $s_{2n}$ and $\tilde{s}_{2n}$ denote the corresponding even moments. Then
\begin{gather*}
 U_n:={s_{2n}\over b_0^2b_1^2\cdots b_{n-1}^2} \le \tilde{U}_n:={\tilde{s}_{2n}\over \tilde{b}_0^2\tilde{b}_1^2\cdots \tilde{b}_{n-1}^2}.
\end{gather*}
\end{prop}

For $n_0\in\mathbb N$ the shifted recurrence sequence $b_n^{(n_0)}:=b_{n+n_0}$, $n\ge 0$ leads to moments $s_n^{(n_0)}$ and
\begin{gather*}
U_n^{(n_0)}:=\frac{s_{2n}^{(n_0)}}{\big(b_0^{(n_0)}\cdots b_{n-1}^{(n_0)}\big)^2}=\frac{s_{2n}^{(n_0)}}{b_{n_0}^2\cdots b_{n_0+n-1}^2}.
\end{gather*}

\begin{prop}\label{thm:shiftU} The expressions above satisfy
\begin{gather}\label{eq:shiftU1}
s_{2n+2}\ge b_0^2 s_{2n}^{(1)},\qquad U_{n+1}\ge U_n^{(1)}.
\end{gather}
If $b_n\le b_n^{(n_0)}$ for $n\ge 0$ and $\lim b_n^{1/n}=1$, then
\begin{gather}\label{eq:shiftU2}
\limsup_{n\to\infty} U_n^{1/n}=\limsup_{n\to\infty}\big(U_n^{(n_0)}\big)^{1/n}.
\end{gather}
The latter holds for $b_n=\lambda (n+1)^c$, $\lambda, c>0$.
\end{prop}

\begin{proof} Define
\begin{gather*}
J_{(1)}=\begin{pmatrix} 0 & 0 & 0 & 0 & \cdots\\
 0 & 0 & b_1 & 0 & \cdots\\
 0 & b_1 & 0 & b_2 & \cdots\\
 0 & 0 & b_2 & 0 & \cdots\\
 \vdots & \vdots & \vdots & \vdots & \ddots
\end{pmatrix}=
\begin{pmatrix} 0 & 0\\
0 & J^{(1)}
\end{pmatrix},
\end{gather*}
where $J^{(1)}$ is the Jacobi matrix \eqref{eq:Jacobimat} corresponding to $b_n^{(1)}$ and the matrix to the right is a block matrix. The upper left $0$ is a scalar and the upper right $0$ is an infinite row vector of zeros. In block matrix notation we then have
\begin{gather*}
\big(J_{(1)}\big)^n=\begin{pmatrix} 0 & 0\\
0 & \big(J^{(1)}\big)^n
\end{pmatrix},
\end{gather*}
hence
\begin{gather*}
s_{2n}^{(1)}=\big((J_{(1)})^{2n} \delta_1,\delta_1\big).
\end{gather*}
This gives
\begin{gather*}
s_{2n+2} = \big(J^{2n+2}\delta_0,\delta_0\big)=\big(J^{2n}J\delta_0,J\delta_0\big)=b_0^2\big(J^{2n}\delta_1,\delta_1\big)
\ge b_0^2 \big((J_{(1)})^{2n} \delta_1,\delta_1\big)=b_0^2s_{2n}^{(1)},
\end{gather*}
hence $U_{n+1}\ge U_n^{(1)}$.

If $b_n\le b_n^{(n_0)}$ we have $s_{2n}\le s_{2n}^{(n_0)}$ by Proposition~\ref{thm:increasofbn}, and therefore for $n>n_0$
\begin{gather*}
U_n\le \frac{b_{n_0}^2b_{n_0+1}^2\cdots b_{n_0+n-1}^2}{b_0^2 b_1^2\cdots b_{n-1}^2}U_n^{(n_0)}\le \frac{b_n^2b_{n+1}^2\cdots b_{n+n_0-1}^2}{b_0^2 b_1^2\cdots b_{n_0-1}^2}U_{n+n_0}.
\end{gather*}
If $\lim b_n^{1/n}=1$ we get \eqref{eq:shiftU2}.
\end{proof}

\begin{lem}\label{thm:b_i/b_{i+1}} Assume $(b_n)$ satisfies
\begin{gather*}
\sup_{i,j\ge 0}\frac{b_i}{b_{i+j}}=M<\infty.
\end{gather*}
Then
\begin{gather}\label{eq:Un+1}
U_{n+1}\le 4M^2U_n,\qquad s_{2n+2}\le 4M^2b_n^2s_{2n}.
\end{gather}
\end{lem}

\begin{proof} By the recurrence equation \eqref{eq:u} we get
\begin{gather*}
u_{n+1,k}\le M(u_{n,k}+u_{n,k-1}),
\end{gather*}
hence
\begin{gather*}
u_{n+1,k}^2\le 2M^2\big(u_{n,k}^2+u_{n,k-1}^2\big).
\end{gather*}
This gives
\begin{gather*}
U_{n+1}=\sum_{k=0}^{[(n+1)/2]} u_{n+1,k}^2\le 4M^2 U_n,
\end{gather*}
hence $s_{2n+2}\le 4M^2b_n^2s_{2n}$.
\end{proof}

In the following we shall use the notation from \cite{G:R} about $q$-shifted factorials. For $z,q\in\mathbb C$ we write
\begin{gather}\label{eq:qshift}
(z;q)_n=\prod_{j=0}^{n-1}\big(1-zq^{j}\big),\qquad n=1,2,\ldots,\qquad (z;q)_0=1.
\end{gather}

When $|q|<1$ we define the absolutely convergent infinite product
\begin{gather}\label{eq:q-infty}
(z;q)_{\infty}=\prod_{j=0}^\infty\big(1-zq^j\big).
\end{gather}

